\chapter{Acheter de la connectivité}\label{ch:nw:buying-connectivity}

L'\omterm{def:nw:buying-connectivity:sla}{accord de niveau de service} (SLA) d'un fournisseur de cloud pour ses connexions privées promet 99.99~\% de chaque mois aux clients qui achètent deux
connexions sur deux sites, et un crédit de 10~\% des redevances de port du mois s'il ne livre pas. Pour une seule connexion, il promet 95~\%, et
ne verse rien tant que le mois n'a pas perdu plus d'un jour et demi. Les sociétés de trading achètent des circuits selon le même type de conditions~: le fournisseur
promet un pourcentage, verse une fraction de la redevance mensuelle lorsqu'il manque à son engagement, et la perte causée par la panne (les positions qui n'ont pas pu être couvertes, les cotations
qui n'ont pas pu être retirées) reste à la charge de la société. Ce chapitre porte sur ce qu'une société achète, auprès de qui, à quelles conditions, et pourquoi
le second circuit ne vaut la peine d'être acheté que s'il ne partage pas une tranchée avec le premier.

Les chapitres~9 à~14 ont décrit les bâtiments, les cartes et les tracés. Ici, ils deviennent des achats~: \omterm{def:nw:colocation-products-and-how-they-are-sold:mmr}{interconnexions} (cross-connects) et circuits dans un
bâtiment, \omterm{def:nw:buying-connectivity:metro}{circuits métropolitains} entre bâtiments, fibre longue distance et faisceaux hertziens entre villes, et flux fournis par des tiers. Les prix qui sont
publics proviennent des grilles tarifaires des plateformes~; les autres sont négociés, et le chapitre le précise.

\section{Ce qu'on achète~: circuits, longueurs d'onde, bande passante et flux}

\begin{definition}[Circuit métropolitain]\label{def:nw:buying-connectivity:metro}
Un \emph{circuit métropolitain}\index{circuit métropolitain} (metro circuit) est une connexion point à point entre deux bâtiments de la même zone métropolitaine (deux centres de données,
un centre de données et un bureau), vendue comme un \omterm{def:nw:long-haul-fibre:dark}{service allumé} (lit service) d'un débit donné ou construite sur de la fibre louée ou possédée.
\end{definition}

La connectivité d'une société est un empilement d'achats. Dans chaque bâtiment, elle achète des \omterm{def:nw:colocation-products-and-how-they-are-sold:mmr}{interconnexions} (chapitre~9) et les ports de la plateforme~; entre les
bâtiments d'une zone métropolitaine, des \omterm{def:nw:buying-connectivity:metro}{circuits métropolitains}~; entre les villes, des longueurs d'onde ou de la \omterm{def:nw:long-haul-fibre:dark}{fibre noire} (chapitre~13) et de la
bande passante radio (chapitre~14)~; et, pour les données qu'elle ne collecte pas elle-même, des flux livrés par un extranet ou un fournisseur. Une société du New Jersey qui
négocie au NYSE, au Nasdaq et chez Cboe depuis un seul bâtiment achète au moins deux \omterm{def:nw:buying-connectivity:metro}{circuits métropolitains} et les ports des plateformes à l'autre extrémité~; celle qui
négocie contre eux les contrats à terme de Chicago ajoute une route longue distance et sa sauvegarde.

\begin{omfigure}
\begin{tikzpicture}[font=\footnotesize,b/.style={draw,rounded corners=2pt,minimum height=0.6cm,align=center,inner sep=3pt},>=Stealth]
\node[b,fill=omDef!15,minimum width=2.4cm] (f) at (0,0) {baie de la société\\ (Carteret)};
\node[b,fill=omThm!15,minimum width=2.2cm] (v1) at (0,2.0) {ports Nasdaq\\ (même bâtiment)};
\node[b,fill=omThm!15,minimum width=2.2cm] (v2) at (6.4,1.2) {groupe NYSE\\ (Mahwah)};
\node[b,fill=omThm!15,minimum width=2.2cm] (v3) at (6.4,-0.8) {Cboe\\ (Secaucus NY5)};
\node[b,fill=omThm!15,minimum width=2.2cm] (v4) at (11.6,1.2) {CME Globex\\ (Aurora)};
\node[b,fill=gray!10,minimum width=2.2cm] (x) at (0,-2.0) {extranet ou\\ flux du fournisseur};
\draw[->,thick] (f) -- node[right,font=\scriptsize]{interconnexion} (v1);
\draw[->,thick] (f) -- node[sloped,above,font=\scriptsize]{circuit métropolitain} (v2);
\draw[->,thick] (f) -- node[sloped,below,font=\scriptsize]{circuit métropolitain} (v3);
\draw[->,thick,omMeth] (v2.east) -- node[above,font=\scriptsize,omMeth,align=center]{longue distance~:\\ radio, fibre} (v4.west);
\draw[->,thick,gray] (f) -- node[right,font=\scriptsize]{circuit} (x);
\end{tikzpicture}

\omcaption{Connectivité d'une société du New Jersey sous forme d'achats, de manière schématique~: \omterm{def:nw:colocation-products-and-how-they-are-sold:mmr}{interconnexions} et ports dans son propre bâtiment, \omterm{def:nw:buying-connectivity:metro}{circuits métropolitains} vers les autres
deux bâtiments du triangle (chapitre~10), une route longue distance avec sa sauvegarde vers Chicago (chapitres~13 et~14), et un circuit vers un fournisseur pour les
données qu'elle ne collecte pas elle-même.}\label{fig:nw:buying-connectivity:stack}
\end{omfigure}

\begin{dated}{2026-09}{Prix publiés de la connectivité~: ordres de grandeur}\label{dat:nw:buying-connectivity:prices}
La grille tarifaire de la connectivité du groupe NYSE (dernière mise à jour en avril 2025) facture, par mois~: 2\,500~USD pour un circuit IP de 1~gigabit, 11\,000
et 18\,000~USD pour des connexions réseau IP et NMS de 10 et de 40~gigabits, 600~USD pour une \omterm{def:nw:colocation-products-and-how-they-are-sold:mmr}{interconnexion} fibre dans le centre de données, 9\,000 à
44\,000~USD pour de 10 à 200~mégabits de bande passante sans fil entre Mahwah et Secaucus (10\,000 à 45\,000 vers Carteret), et 6\,000~USD pour une
connexion sans fil transportant les données du groupe CME~; des frais non récurrents de 500 à 15\,000~USD s'appliquent. MIAX Pearl facture 15\,000~USD par mois pour une
connexion de 10~gigabits à latence ultra-faible (chapitre~9)~; Deutsche Börse, de 6\,000 à 8\,400~EUR par mois pour ses connexions de co-location de 10~gigabits
(chapitre~11).
\end{dated}

\begin{table}[ht]
\centering\small
\begin{tabular}{lrrl}
\toprule
Poste & Mensuel & Non récurrent & Source \\
 & (USD) & (USD) & \\
\midrule
NYSE réseau IP, circuit de 1~Gb & 2\,500 & 2\,500 & grille tarifaire \\
NYSE réseaux IP et NMS, 10~Gb & 11\,000 & 10\,000 & grille tarifaire \\
NYSE réseaux IP et NMS, 40~Gb & 18\,000 & 10\,000 & grille tarifaire \\
NYSE \omterm{def:nw:colocation-products-and-how-they-are-sold:mmr}{interconnexion} fibre dans le centre de données & 600 & 500 & grille tarifaire \\
NYSE sans fil Mahwah--Secaucus, 10~Mb & 9\,000 & 10\,000 & grille tarifaire \\
NYSE sans fil Mahwah--Secaucus, 200~Mb & 44\,000 & 10\,000 & grille tarifaire \\
NYSE connexion sans fil des données CME & 6\,000 & 5\,000 & grille tarifaire \\
MIAX Pearl 10~Gb à latence ultra-faible & 15\,000 & 0 & dépôt tarifaire \\
\bottomrule
\end{tabular}
\caption{Le catalogue de \texttt{firm.slamodel}~: seuls les prix publiés, chacun avec sa ligne du registre et sa date. La bande passante sans fil coûte à peu près autant
par mégabit et par mois qu'une connexion par fibre de 10~gigabits coûte par gigabit. Données~: \texttt{nw\_buy.price\_rows()}.}\label{tab:nw:buying-connectivity:prices}
\end{table}

Le tableau donne l'échelle du marché. Un mégabit de faisceau hertzien entre deux bâtiments du New Jersey coûte environ 900~USD par mois à 10~mégabits et
220~USD à 200~; un gigabit de fibre entre une baie et le réseau de la plateforme coûte environ 1\,100~USD par mois à 10~gigabits. Ce que vend la radio,
c'est la vitesse, non la bande passante.

\section{Extranets financiers et infrastructure gérée}

\begin{definition}[Extranet financier]\label{def:nw:buying-connectivity:extranet}
Un \emph{extranet financier}\index{extranet financier} (financial extranet) est un réseau privé, exploité par un fournisseur relié à de nombreuses plateformes et sociétés, par lequel
ses clients accèdent aux données de marché et à la saisie d'ordres des plateformes, et communiquent entre eux, sans construire leurs propres circuits vers chacune~; les plateformes admettent ces fournisseurs
dans le cadre d'un accord et leur facturent l'accès qu'ils revendent.
\end{definition}

La définition du Nasdaq montre ce dispositif du point de vue de la plateforme~: un fournisseur d'extranet signe le Nasdaq extranet provider agreement, dispose de
sa propre connexion au Nasdaq, est hébergé en co-location dans une installation du Nasdaq, et donne à ses propres clients accès aux services du Nasdaq, dans le bâtiment ou ailleurs.
Pour une société, l'extranet est un compromis~: une seule connexion vers le fournisseur remplace de nombreuses connexions vers les plateformes, au prix de la latence du fournisseur et du
partage de son réseau avec ses autres clients.

\begin{definition}[Infrastructure gérée]\label{def:nw:buying-connectivity:managed}
L'\emph{infrastructure gérée}\index{infrastructure gérée} (managed infrastructure) est l'équipement et la connectivité qu'un fournisseur installe, possède ou loue, et exploite
pour le compte d'une société (serveurs, commutateurs, circuits, gestionnaires de données de marché), en général dans les baies du fournisseur près des plateformes, et vendue comme
un service mensuel.
\end{definition}

L'\omterm{def:nw:buying-connectivity:managed}{infrastructure gérée} relève de la question construire ou acheter du livre~16 (construire contre acheter), à l'échelle d'une baie~: la société qui l'achète obtient
une présence sur une plateforme en quelques semaines au lieu de quelques mois, et confie au fournisseur le contrôle de la latence de chaque élément du chemin. Les sociétés dont
l'edge tient à la latence achètent l'espace du bâtiment et construisent le reste~; celles dont l'edge est ailleurs achètent aussi le reste.

\section{Fournisseurs de bande passante sans fil}

Les routes radio se vendent de trois façons. Les plateformes vendent de la bande passante ou des données sur leurs propres liaisons~: ICE entre son centre de données de Mahwah et les autres bâtiments du New
Jersey (avec Anova Financial Networks) et depuis Toronto~; SIX entre Zurich et Francfort~; Euronext de Londres à Bergame avec
McKay Brothers. Les opérateurs de réseau vendent de la bande passante sur les leurs. Et les fournisseurs de données vendent un flux de données de marché livré par radio, comme Quincy Data le faisait
depuis Aurora au chapitre~10. Ce qu'ils ont en commun, c'est la rareté~: des mégabits, non des gigabits~; un ensemble fixe de points d'extrémité~; et une sauvegarde en fibre que
la société doit aussi acheter, ou accepter de se passer lorsqu'il pleut.

\section{Contrats et niveaux de service}

\begin{definition}[Accord de niveau de service, crédit de service]\label{def:nw:buying-connectivity:sla}
Un \emph{accord de niveau de service}\index{accord de niveau de service} (SLA, service-level agreement) est la partie d'un contrat de fourniture qui énonce les performances auxquelles le fournisseur
s'engage (disponibilité, latence, délai de réparation), la manière dont elles sont mesurées et ce que le fournisseur doit lorsqu'il manque à ses engagements. Un
\emph{crédit de service}\index{crédit de service} est ce qu'il doit~: en général un pourcentage des frais mensuels du service défaillant, imputé sur les factures suivantes.
\end{definition}

\begin{definition}[Temps moyen entre pannes, temps moyen de réparation]\label{def:nw:buying-connectivity:mtbf}
Le \emph{temps moyen entre pannes}\index{temps moyen entre pannes} (MTBF, mean time between failures) d'un composant est la durée moyenne pendant laquelle il fonctionne entre deux pannes~; son
\emph{temps moyen de réparation}\index{temps moyen de réparation} (MTTR, mean time to repair) est la durée moyenne qui sépare une panne de son rétablissement. Sa disponibilité en régime permanent est
$\mathrm{MTBF}/(\mathrm{MTBF}+\mathrm{MTTR})$.
\end{definition}

\begin{dated}{2026-09}{Un barème de crédits publié}\label{dat:nw:buying-connectivity:sla}
L'\omterm{def:nw:buying-connectivity:sla}{accord de niveau de service} d'AWS pour Direct Connect (dernière mise à jour en juillet 2026) comporte trois engagements~: une disponibilité mensuelle de 99.99~\% pour les déploiements redondants sur plusieurs sites,
de 99.9~\% pour ceux sur plusieurs sites sans redondance, et de 95.0~\% pour une connexion unique. Leur non-respect ouvre droit à un crédit de 10~\%, 25~\% ou
100~\% des frais d'accès au port des connexions concernées, aux seuils respectifs de (99.99, 99.0, 95.0), (99.9, 99.0, 95.0) et (95.0, 92.5, 90.0)~\%.
Les crédits constituent le seul recours, et les connexions uniques ne sont pas recommandées en production.
\end{dated}

\begin{definition}[Diversité de tracé]\label{def:nw:buying-connectivity:diversity}
La \emph{diversité de tracé}\index{diversité de tracé} est la propriété de deux circuits ou plus entre les mêmes points d'extrémité de ne partager aucun élément physique
dont la panne arrêterait les deux~: ni conduit, ni câble, ni entrée de bâtiment, ni salle, ni équipement, ni alimentation électrique.
\end{definition}

\begin{proposition}[Deux circuits]\label{prop:nw:buying-connectivity:two}
Supposons que chacun de deux circuits soit indisponible une fraction $u$ du temps, de façon indépendante, et que les événements de mode commun qui arrêtent les deux surviennent une fraction
$c$ du temps. La paire est alors indisponible une fraction $u^2 + c - u^2 c \approx u^2 + c$. Lorsque $c$ est de l'ordre de $u$, le second circuit
n'apporte presque rien~; seule la diversité, $c \to 0$, donne $u^2$.
\end{proposition}

\begin{proof}
La paire est en panne lorsque les deux le sont indépendamment (probabilité $u^2$) ou qu'un événement de mode commun est en cours ($c$)~; la formule du crible
donne l'union.
\end{proof}

\omcode{code/firm/slamodel/firm_slamodel.py}{78}{84}{Disponibilité en régime permanent d'un montage~: tous les circuits en panne en même temps, ou un événement de mode commun.}\label{lst:nw:buying-connectivity:avail}

\begin{omfigure}
\begin{tikzpicture}[font=\footnotesize,b/.style={draw,rounded corners=2pt,minimum height=0.55cm,align=center,inner sep=2pt},>=Stealth]
\foreach \y/\lab in {2.4/{un circuit},0.9/{deux circuits dans une même conduite},-0.8/{circuits diversifiés}}{
  \node[b,fill=omDef!15,minimum width=1.6cm] at (0,\y) {société};
  \node[b,fill=omThm!15,minimum width=1.6cm] at (9,\y) {plateforme};
  \node[anchor=west,font=\scriptsize] at (-1,\y+0.55) {\lab};}
\draw[thick] (0.8,2.4) -- (8.2,2.4);
\draw[gray,line width=4pt,opacity=0.3] (2.2,0.9) -- (6.8,0.9);
\draw[thick] (0.8,1.0) -- (8.2,1.0); \draw[thick] (0.8,0.8) -- (8.2,0.8);
\node[gray,font=\scriptsize] at (4.5,0.5) {une conduite~: une seule pelleteuse coupe les deux};
\draw[thick] (0.8,-0.7) -- (2,-0.2) -- (7,-0.2) -- (8.2,-0.7);
\draw[thick] (0.8,-0.9) -- (2,-1.4) -- (7,-1.4) -- (8.2,-0.9);
\node[font=\scriptsize] at (4.5,0.05) {tracé A}; \node[font=\scriptsize] at (4.5,-1.65) {tracé B~: une autre rue, une autre entrée, un autre fournisseur};
\end{tikzpicture}

\omcaption{Trois montages, schématiquement~: un circuit~; deux circuits dont les fibres partagent une conduite, de sorte qu'un seul événement (une coupure, un incendie) arrête les deux~; et deux
circuits diversifiés, qui ne partagent rien entre les salles de la société et celles de la plateforme.}\label{fig:nw:buying-connectivity:designs}
\end{omfigure}

Le modèle attribue à chaque circuit deux pannes par an, avec une réparation moyenne de quatre heures (MTBF \qty{4383}{\hour}, MTTR \qty{4}{\hour}), et
la conduite commune subit une demi-panne de mode commun par an, avec une réparation moyenne de douze heures (une fibre coupée est épissée, non remplacée)~: ce sont des hypothèses, signalées comme telles.
Un circuit est alors indisponible 0.091~\% de l'année, soit 480 minutes. Deux circuits dans une même conduite sont indisponibles 360 minutes~: les six heures annuelles de la conduite, presque en totalité,
puisque deux pannes indépendantes coïncident moins d'une demi-minute par an. Deux circuits diversifiés sont indisponibles 0.4 minute par an.

\begin{table}[ht]
\centering\small
\begin{tabular}{@{}>{\raggedright\arraybackslash}p{2.8cm}>{\centering\arraybackslash}p{2.1cm}>{\centering\arraybackslash}p{2.1cm}>{\centering\arraybackslash}p{2.1cm}>{\centering\arraybackslash}p{2.1cm}>{\centering\arraybackslash}p{2.1cm}@{}}
\toprule
Montage & Disponibilité & Espérance & 95\up{e}~centile & Coût annuel & Perte à 2\,000 \\
 & (\%) & (min/an) & (min/an) & (kUSD) & USD/min (kUSD) \\
\midrule
Un circuit & 99.9088 & 479.6 & 1\,541 & 139.2 & 959 \\
Deux circuits dans une même conduite & 99.9315 & 360.4 & 1\,992 & 278.4 & 721 \\
Deux circuits diversifiés & 99.99992 & 0.4 & 0 & 320.2 & 0.9 \\
\bottomrule
\end{tabular}
\caption{Simulation de trois montages de connexions NYSE à 10~gigabits (prix publiés~; les taux de panne, les événements de mode commun de la conduite, une
prime de 15~\% pour la diversité et la perte de trading par minute sont des hypothèses)~: disponibilité en régime permanent et minutes d'indisponibilité espérées, le 95\up{e}~centile
d'une année simulée (1\,000 années), coût annuel avec une \omterm{def:nw:colocation-products-and-how-they-are-sold:mmr}{interconnexion} par circuit, et la perte de trading espérée. Données~:
\texttt{nw\_buy.results()}.}\label{tab:nw:buying-connectivity:designs}
\end{table}

\begin{omfigure}
\begin{tikzpicture}
\begin{semilogyaxis}[width=10.5cm,height=5.8cm,xmin=100,xmax=460,ymin=0.1,ymax=2000,xlabel={coût annuel (milliers de USD)},
  ylabel={indisponibilité (min par an, log)},tick label style={font=\footnotesize},label style={font=\small},grid=major,grid style={gray!25},
  table/col sep=comma,nodes near coords,point meta=explicit symbolic,every node near coord/.append style={font=\scriptsize,anchor=west}]
\addplot[only marks,mark=*,mark size=2.6pt,omThm] table[x=cost_k,y=expected_min,meta=design,col sep=comma,row sep=\\]{design,cost_k,expected_min,loss_k\\ un circuit,139.2,479.562,959.12\\ deux dans une conduite,278.4,360.437,720.87\\ deux diversifiés,320.2,0.437,0.87\\};
\end{semilogyaxis}
\end{tikzpicture}

\omcaption{Minutes d'indisponibilité espérées par an en fonction du coût annuel pour les trois montages du \cref{tab:nw:buying-connectivity:designs}~: le second circuit dans
la même conduite double le coût et supprime un quart de l'indisponibilité~; la diversité ajoute 15~\% et supprime presque tout. Données~: \texttt{fig\_buy.py}.}\label{fig:nw:buying-connectivity:cost}
\end{omfigure}

\omcode{code/firm/slamodel/firm_slamodel.py}{118}{129}{Une année simulée~: les pannes de chaque circuit, les événements de mode commun et les minutes pendant
lesquelles tous les circuits étaient en panne en même temps.}\label{lst:nw:buying-connectivity:sim}

\begin{omfigure}
\begin{tikzpicture}
\begin{semilogyaxis}[width=11cm,height=6cm,xmin=0,xmax=100,ymin=0.01,ymax=20000,xlabel={centile des années simulées},
  ylabel={minutes d'indisponibilité dans l'année (log)},tick label style={font=\footnotesize},label style={font=\small},grid=major,grid style={gray!25},
  legend columns=3,legend style={font=\footnotesize,at={(0.5,-0.25)},anchor=north,draw=none},table/col sep=comma]
\addplot[omDef,thick,no markers] table[x=pct,y=minutes]{figdata/networks/15-buying-connectivity/ecdf_single.csv};
\addplot[omMeth,thick,dashed,no markers] table[x=pct,y=minutes]{figdata/networks/15-buying-connectivity/ecdf_duct.csv};
\addplot[omThm,thick,no markers] table[x=pct,y=minutes]{figdata/networks/15-buying-connectivity/ecdf_diverse.csv};
\legend{un circuit,deux circuits dans une même conduite,deux circuits diversifiés}
\end{semilogyaxis}
\end{tikzpicture}

\omcaption{Simulation~: minutes d'indisponibilité dans l'année selon le centile parmi 1\,000 années simulées (zéro représenté à 0.01). Deux circuits dans une même conduite ont davantage
d'années sans panne qu'un circuit seul, et de plus mauvaises années~: un événement sur la conduite dure douze heures en moyenne. Les circuits diversifiés ne tombent presque jamais
en panne ensemble. Données~: \texttt{fig\_buy.py}.}\label{fig:nw:buying-connectivity:ecdf}
\end{omfigure}

Que rapportent ces pannes en \omterm{def:nw:buying-connectivity:sla}{crédits de service}~? En appliquant le barème publié de l'\cref{dat:nw:buying-connectivity:sla} à chaque mois simulé,
le circuit seul, selon les conditions d'une connexion unique, rapporte en moyenne 5.50~USD par an~: une panne de quatre heures ramène le mois à 99.44~\%, bien au-dessus de
95~\%. Les deux circuits dans une même conduite, selon les conditions sans redondance, rapportent environ 2\,560~USD par an~; la paire diversifiée, selon les conditions avec redondance, ne rapporte rien.
Face à des pertes de l'ordre de centaines de milliers de dollars dans le modèle, les crédits ne sont pas une assurance~: ils signalent ce que le fournisseur
s'attend à livrer, et l'entreprise s'assure elle-même par la conception.

\section{Ordres de grandeur des prix}

Les prix publiés fixent un plancher au budget. Une entreprise dotée de deux connexions diversifiées de 10~gigabits vers le groupe NYSE paie environ 320\,000~USD par an
pour elles dans le modèle~; les deux connexions à ultra-faible latence de la même entreprise vers une bourse d'options coûtent 360\,000~USD (chapitre~9)~; une liaison radio de 200
mégabits entre deux immeubles du New Jersey coûte plus d'un demi-million par an. Ce qui n'est pas publié (services d'extranet, \omterm{def:nw:buying-connectivity:metro}{circuits métropolitains} entre
immeubles privés, longueurs d'onde longue distance, baies gérées) est tarifé sur cotation, et le chapitre~29 traite ces lignes comme des estimations assorties de fourchettes, jamais comme
des faits.

\begin{method}[Acheter un circuit]\label{met:nw:buying-connectivity:buy}
\begin{enumerate}
\item Dire à quoi sert le circuit~: quelles stratégies, quelle latence, quelle bande passante, ce qui se passe quand il tombe en panne.
\item Demander à chaque fournisseur le tracé (conduites, rues, entrées d'immeubles, salles), la latence unidirectionnelle de baie à baie et la manière dont elle est mesurée, ainsi que son
  historique de pannes.
\item Acheter de la diversité, pas des doublons~: vérifier que le second circuit ne partage rien avec le premier, jusqu'à l'entrée de l'immeuble et à l'alimentation électrique.
\item Lire l'\omterm{def:nw:buying-connectivity:sla}{accord de niveau de service} (SLA) pour ce qu'il mesure (disponibilité mensuelle, délai de réparation), les crédits et les exclusions~; comparer les crédits avec sa propre perte par
  minute.
\item Inscrire chaque prix, chaque condition et chaque date de renouvellement dans le budget avec sa source (chapitre~29), et revérifier avant le renouvellement.
\end{enumerate}
\end{method}

\section{Tutoriel~: une année de pannes}\label{tut:nw:buying-connectivity:tut}

\begin{tutorial}
\textbf{Objectif.} Chiffrer trois conceptions à partir des tarifs publiés, simuler leurs pannes et comparer les crédits aux pertes.
\textbf{État final~:} la \cref{fig:nw:buying-connectivity:ecdf} et les \cref{tab:nw:buying-connectivity:prices,tab:nw:buying-connectivity:designs}.
\begin{tutsteps}
\item \textbf{Le catalogue.} \texttt{firm\_slamodel.CATALOGUE} regroupe les prix publiés avec leurs lignes sourcées~; \texttt{SCHEDULES} les paliers
  de crédit publiés.
\item \textbf{Les conceptions.} \texttt{nw\_buy.designs()} construit un circuit, deux dans une même conduite et deux diversifiés à partir de \texttt{ASSUME}.
\item \textbf{Régime permanent et simulation.} \texttt{design\_availability} (le \cref{lst:nw:buying-connectivity:avail}) et \texttt{simulate\_year}
  (le \cref{lst:nw:buying-connectivity:sim})~; \texttt{results()} remplit le \cref{tab:nw:buying-connectivity:designs}.
\item \textbf{Crédits.} \texttt{monthly\_credits} découpe chaque année simulée en mois et applique le barème.
\end{tutsteps}
\textbf{Pour aller plus loin.} Rendre les événements de conduite plus rares mais plus longs~; placer les deux circuits diversifiés chez le même fournisseur et ajouter une
panne commune au fournisseur~; remplacer le coût par minute de perte par le chiffre de sa propre stratégie.
\end{tutorial}

\section{Construction~: le modèle de SLA}\label{bld:nw:buying-connectivity:slamodel}

\begin{build}
\bfield{Rôle} Des circuits tels que chiffrés et sourcés dans le catalogue, leur disponibilité seule et dans les conceptions, les crédits qu'un contrat verse et les pertes
qu'il ne couvre pas~: l'entrée des chapitres~28 (résilience) et~29 (le plan).
\bfield{Interface} \texttt{firm\_slamodel}~: \texttt{Tier}, \texttt{Circuit}, \texttt{Design}, \texttt{CATALOGUE}, \texttt{SCHEDULES},
\texttt{availability}, \texttt{design\_availability}, \texttt{simulate\_year}, \texttt{monthly\_down\_min}, \texttt{uptime\_pct}, \texttt{credit},
\texttt{downtime\_cost}, \texttt{expected\_down\_min}, \texttt{nines}.
\bfield{Règles} Les prix ne viennent que du catalogue, avec leur source et leur date~; les taux de panne sont des hypothèses déclarées~; les événements de mode commun d'une conception
sont explicites, jamais nuls par défaut lorsque les circuits partagent quoi que ce soit.
\bfield{Tests d'acceptation} \texttt{code/firm/slamodel/tests/}~: disponibilité calculée à la main pour les trois conceptions, simulation confrontée au régime permanent,
découpage mensuel et paliers de crédit.
\bfield{Pour aller plus loin} Des durées de réparation qui dépendent de la panne (une coupure contre une carte), des fenêtres de maintenance annoncées à l'avance, et un mode commun
propre au fournisseur entre ses clients.
\end{build}

\begin{omsources}
\item AWS Direct Connect service-level agreement (juillet 2026).
\item Nasdaq Data News \#2017-8 (fournisseurs d'extranet)~; NYSE group connectivity fee schedule (avril 2025).
\item Chapitres~9 à~14 pour les descriptions que les plateformes et les fournisseurs donnent de leurs propres connexions.
\end{omsources}

\section{Exercices}

\begin{exercise}[$\star$]\label{exo:nw:buying-connectivity:1}
Un circuit tombe en panne deux fois par an et met quatre heures à être réparé à chaque fois. Calculer sa disponibilité et le nombre de minutes d'indisponibilité dans l'année.
\end{exercise}

\begin{exercise}[$\star$]\label{exo:nw:buying-connectivity:2}
D'après le \cref{tab:nw:buying-connectivity:prices}, calculer le coût mensuel d'un mégabit de bande passante sans fil entre Mahwah et Secaucus à 10 mégabits, puis à 200
mégabits.
\end{exercise}

\begin{exercise}[$\star$]\label{exo:nw:buying-connectivity:3}
Selon le barème à connexion unique de l'\cref{dat:nw:buying-connectivity:sla}, calculer pendant combien de temps une connexion doit être indisponible dans un mois de 30 jours avant
qu'un crédit ne soit dû.
\end{exercise}

\begin{exercise}[$\star\star$]\label{exo:nw:buying-connectivity:4}
En utilisant la \cref{prop:nw:buying-connectivity:two}, expliquer pourquoi deux circuits dans une même conduite apportent si peu par rapport à un seul.
\end{exercise}

\begin{exercise}[$\star\star$]\label{exo:nw:buying-connectivity:5}
Énumérer cinq éléments que deux circuits «~diversifiés~» peuvent encore partager.
\end{exercise}

\begin{exercise}[$\star\star$]\label{exo:nw:buying-connectivity:6}
Dire quand un \omterm{def:nw:buying-connectivity:extranet}{extranet financier} est la bonne manière d'atteindre une plateforme, et quand il ne l'est pas.
\end{exercise}

\begin{exercise}[$\star\star\star$]\label{exo:nw:buying-connectivity:7}
\emph{Programmation.} Avec \texttt{firm.slamodel}, trouver le taux de mode commun à partir duquel deux circuits dans une même conduite ne valent pas mieux qu'un circuit seul.
\end{exercise}

\begin{exercise}[$\star\star\star$]\label{exo:nw:buying-connectivity:8}
\emph{Trouver la faille.} «~Notre fournisseur garantit 99.99~\% et verse des crédits, donc une panne ne nous coûte rien.~»
\end{exercise}

\section{Problème~: deux circuits ne valent pas deux fois mieux}

\begin{problem}[{Problème du week-end --- la redondance qui n'en est pas une}]\label{pb:nw:buying-connectivity:1}
Une entreprise se connecte au groupe NYSE par des liaisons de 10~gigabits. Chacune tombe en panne deux fois par an et met quatre heures en moyenne à être réparée. Ses
deux circuits passent dans une même conduite, qui subit un demi-événement de mode commun par an avec douze heures de réparation~; une paire diversifiée coûterait 15~\% de plus. L'entreprise
perd 2\,000~USD pour chaque minute sans connexion.

\medskip\noindent\textbf{Partie I --- Un circuit.}
\begin{enumerate}
\item Quelles sont la disponibilité du circuit et les minutes d'indisponibilité attendues dans l'année~?
\item Combien coûte-t-il par an avec son \omterm{def:nw:colocation-products-and-how-they-are-sold:mmr}{interconnexion}~?
\item Quelle est la perte annuelle attendue~?
\item Quel crédit le barème à connexion unique verse-t-il pour une panne de quatre heures~?
\end{enumerate}

\medskip\noindent\textbf{Part II --- Two circuits in one duct.}
\begin{enumerate}[resume]
\item Quelles sont la disponibilité de la paire et les minutes d'indisponibilité attendues~?
\item Quelle part de ce total revient au conduit~?
\item Combien coûte la paire, et quelle perte évite-t-elle par rapport à un seul circuit~?
\item À quoi ressemble une mauvaise année (95\up{e}~centile)~?
\end{enumerate}

\medskip\noindent\textbf{Partie III --- Deux circuits diversifiés.}
\begin{enumerate}[resume]
\item Quelles sont la disponibilité de la paire et les minutes d'indisponibilité attendues~?
\item Combien coûte la \omterm{def:nw:buying-connectivity:diversity}{diversité de tracé} par rapport au conduit, et que fait-elle économiser~?
\item Quels crédits les trois architectures rapportent-elles en moyenne par an~?
\item Qu'est-ce qui ferait tomber la paire diversifiée en même temps malgré tout~?
\end{enumerate}

\medskip\noindent\textbf{Partie IV --- Le verdict.}
\begin{enumerate}[resume]
\item Énoncer le \emph{résultat nommé}~: la disponibilité de deux circuits partageant un conduit comparée à celle de deux circuits diversifiés, leurs minutes d'indisponibilité annuelles attendues,
  et le \omterm{def:nw:buying-connectivity:sla}{crédit de service} face à la perte de trading.
\item Quelle hypothèse modifie le plus la réponse~?
\item Comment vérifieriez-vous l'affirmation d'un fournisseur sur la diversité de ses tracés~?
\item Que doit mesurer le SLA pour une société de trading, au-delà de la disponibilité mensuelle~?
\item Comment les systèmes de la société utilisent-ils deux circuits~: en actif et en secours, ou en actif sur les deux~?
\item Où cette analyse s'inscrit-elle dans le plan de connectivité du chapitre~29~?
\item Qu'apporte un second circuit s'il partage le conduit~?
\item En une phrase~: qu'est-ce que la redondance~?
\end{enumerate}
\end{problem}

\section{Questions d'entretien}

\begin{interviewq}[$\star$ \iqroles{développeur}]\label{iq:nw:buying-connectivity:1}
Qu'est-ce qu'un SLA, et que compense un \omterm{def:nw:buying-connectivity:sla}{crédit de service}~?
\end{interviewq}

\begin{interviewq}[$\star\star$ \iqroles{développeur}]\label{iq:nw:buying-connectivity:2}
Vous disposez de deux circuits vers la même plateforme auprès de deux fournisseurs. Comment peuvent-ils quand même tomber en panne ensemble~?
\end{interviewq}

\begin{interviewq}[$\star\star$ \iqroles{développeur, trader}]\label{iq:nw:buying-connectivity:3}
Quand utiliseriez-vous un \omterm{def:nw:buying-connectivity:extranet}{extranet financier} plutôt qu'une connexion directe à une plateforme~?
\end{interviewq}

\begin{interviewq}[$\star\star$ \iqroles{développeur}]\label{iq:nw:buying-connectivity:4}
Comment le MTBF et le MTTR se combinent-ils en disponibilité, et que fait un second circuit indépendant à celle-ci~?
\end{interviewq}

\begin{interviewq}[$\star\star$ \iqroles{trader}]\label{iq:nw:buying-connectivity:5}
La bande passante radio coûte des centaines de fois plus cher par mégabit que la fibre. Pourquoi les sociétés l'achètent-elles~?
\end{interviewq}

\begin{interviewq}[$\star\star\star$ \iqroles{développeur, chercheur}]\label{iq:nw:buying-connectivity:6}
Concevez la connectivité d'un nouveau bureau de trading vers trois plateformes dans deux villes, pour une disponibilité de 99.999~\%. Qu'achetez-vous et que vérifiez-vous~?
\end{interviewq}
